%% neuralmi_methods_paper.tex
%% A detailed methods paper draft for NeuralMI.
%% Author:  Eslam Abdelaleem, Georgia Institute of Technology
%% Status:  DRAFT — intended as a complete self-contained starting point
%%          for a final manuscript.  All section content is drawn directly
%%          from the library source code and documentation.
%%
%% Build:   pdflatex neuralmi_methods_paper.tex
%%          bibtex   neuralmi_methods_paper
%%          pdflatex neuralmi_methods_paper.tex  (x2)
%%
%% Required packages: amsmath, amssymb, mathtools, algorithm2e, listings,
%%   booktabs, hyperref, cleveref, natbib, graphicx, xcolor, geometry,
%%   microtype, setspace.

\documentclass[12pt]{article}

%% ---- Packages ---------------------------------------------------------------
\usepackage[margin=1in]{geometry}
\usepackage{amsmath,amssymb,mathtools}
\usepackage{bm}              % bold math
\usepackage{graphicx}
\usepackage{xcolor}
\usepackage{booktabs}
\usepackage{array}
\usepackage{multirow}
\usepackage{hyperref}
\usepackage{cleveref}
\usepackage{microtype}
\usepackage{setspace}
\usepackage{natbib}
\usepackage{enumitem}
\usepackage{caption}
\usepackage{subcaption}
\usepackage[ruled,vlined,linesnumbered]{algorithm2e}
\usepackage{listings}
\usepackage{parskip}

%% ---- Code listing style -----------------------------------------------------
\definecolor{codegray}{gray}{0.95}
\definecolor{codecomment}{rgb}{0.4,0.4,0.4}
\definecolor{codestring}{rgb}{0.13,0.47,0.13}
\definecolor{codekeyword}{rgb}{0.0,0.0,0.8}
\lstset{
  language=Python,
  basicstyle=\ttfamily\small,
  backgroundcolor=\color{codegray},
  commentstyle=\color{codecomment},
  stringstyle=\color{codestring},
  keywordstyle=\color{codekeyword},
  showstringspaces=false,
  breaklines=true,
  frame=single,
  framesep=4pt,
  xleftmargin=8pt,
  xrightmargin=8pt,
  numbers=none,
  captionpos=b,
}

%% ---- Macros -----------------------------------------------------------------
\newcommand{\nmi}{\texttt{NeuralMI}}
\newcommand{\api}[1]{\texttt{#1}}
\newcommand{\MI}[1]{I(#1)}
\newcommand{\CMI}[2]{I(#1\,|\,#2)}
\newcommand{\E}[1]{\mathbb{E}\!\left[#1\right]}
\newcommand{\KL}[2]{D_{\mathrm{KL}}\!\left(#1\,\|\,#2\right)}
\newcommand{\R}{\mathbb{R}}

\hypersetup{
  colorlinks=true, linkcolor=blue, citecolor=blue, urlcolor=blue,
  pdftitle={NeuralMI: A Python Toolbox for Rigorous Mutual Information
            Estimation in Neuroscience},
  pdfauthor={Eslam Abdelaleem},
}

%% ---- Title block ------------------------------------------------------------
\title{%
  \textbf{NeuralMI: A Python Toolbox for Rigorous Neural\\
  Mutual Information Estimation in Neuroscience}%
}
\author{%
  Eslam Abdelaleem\\
  \small School of Physics, Georgia Institute of Technology\\
  \small \href{mailto:eabdelaleem3@gatech.edu}{eabdelaleem3@gatech.edu}
}
\date{March 2026}

%% ============================================================================
\begin{document}
\maketitle
\onehalfspacing

%% ---- Abstract ---------------------------------------------------------------
\begin{abstract}
We introduce \nmi{}, a Python library that provides neuroscientists with a
complete, end-to-end workflow for estimating mutual information (MI) from
complex neural data.  MI quantifies how much information one variable carries
about another without assuming any particular form of dependence, making it
uniquely suited to the high-dimensional, nonlinear signals typical of modern
systems neuroscience.  The library is built around neural network–based
estimators---InfoNCE and SMILE---that scale to high-dimensional continuous
recordings, spike trains, and categorical state sequences.  All functionality
is exposed through a single entry point, \api{nmi.run()}, which dispatches to
one of nine analysis modes depending on the scientific question: a point
estimate (\api{estimate}), parallelized hyperparameter exploration (\api{sweep}),
temporal lag analysis (\api{lag}), spike-timing precision thresholds
(\api{precision}), bias-corrected estimation with confidence intervals
(\api{rigorous}), latent dimensionality of a neural population (\api{dimensionality}),
conditional MI (\api{conditional}), transfer entropy (\api{transfer}), and
all-to-all pairwise connectivity matrices (\api{pairwise}).  The \api{rigorous}
mode implements a subsampling-and-linear-extrapolation bias-correction procedure
that exploits the theoretically predicted $O(1/N)$ finite-sample bias scaling.
The \api{dimensionality} mode implements a spectral approach based on the
cross-covariance Participation Ratio of learned embeddings, avoiding the
computationally expensive bottleneck sweep used in prior work.  \nmi{} is
freely available at \url{https://github.com/eslam-abdelaleem/NeuralMI} under
the MIT licence.
\end{abstract}

\tableofcontents
\newpage

%% ============================================================================
\section{Introduction}
\label{sec:intro}

Mutual information (MI) is a foundational quantity in information theory that
measures the statistical dependence between two random variables $X$ and $Y$,
capturing \emph{all} forms of dependence—linear and nonlinear—in a single
scalar.  In neuroscience, MI has been applied to quantify the information
encoded by neural populations about sensory stimuli or motor commands, to
measure the functional connectivity between brain regions, and to characterise
the temporal precision of neural codes
\citep{borst1999information,quian2009measuring,timme2018tutorial}.

Despite its theoretical appeal, estimating MI from finite neural recordings is
difficult.  For low-dimensional data, non-parametric estimators such as the
$k$-nearest-neighbour (kNN) method of \citet{kraskov2004estimating} provide
reliable estimates.  However, kNN estimators scale poorly with dimension:
their sample complexity grows exponentially in the intrinsic dimension of the
data, and modern multi-electrode or EEG recordings often have dozens to
hundreds of channels.  Neural network–based estimators overcome the
dimensionality barrier by learning a lower-dimensional representation, but
introduce a different problem: with the finite sample sizes typical of
neuroscience experiments, they systematically \emph{overestimate} the true MI.
This upward bias can be substantial---sometimes exceeding the true MI
itself---yet existing implementations of neural estimators provide only bare
point estimates, with no correction mechanism and no confidence intervals.

\nmi{} is designed to close this gap.  Its key contributions are:

\begin{enumerate}[leftmargin=*]
  \item \textbf{Automated bias correction.}  The \api{rigorous} mode
        implements the subsampling-and-extrapolation procedure introduced in
        \citet{abdelaleem2025accurate}, which exploits the theoretical
        $O(1/N)$ bias scaling to obtain a debiased MI estimate with a
        confidence interval.

  \item \textbf{Nine analysis modes under a single API.}  From a quick point
        estimate to transfer entropy and all-to-all connectivity matrices,
        every analysis is triggered with the same call,
        \api{nmi.run(\ldots, mode=\ldots)}.

  \item \textbf{Neuroscience-native data handling.}  Three processor
        classes—Continuous, Spike, and Categorical—convert raw recordings into
        the unified tensor format that all models expect, with
        memory-efficient non-overlapping windowing and epoch-level jitter.

  \item \textbf{Principled data splitting.}  Two strategies prevent
        train–test leakage: random shuffling for IID data, and
        \emph{blocked} splitting for temporally autocorrelated recordings
        (the safer default).

  \item \textbf{Efficient dimensionality estimation.}  The \api{dimensionality}
        mode computes the Participation Ratio (PR) of the singular-value
        spectrum of the cross-covariance of learned embeddings, replacing
        the computationally expensive bottleneck sweep of prior work
        \citep{gulati2026mutual}.
\end{enumerate}

The remainder of this paper is organized as follows.
\Cref{sec:theory} develops the theoretical background.
\Cref{sec:architecture} describes the library architecture and neural
network components.
\Cref{sec:data} covers data handling and processors.
\Cref{sec:modes} provides a detailed methods description of each of the
nine analysis modes.
\Cref{sec:generators} introduces the built-in synthetic data generators.
\Cref{sec:results} describes the \api{Results} object and utility functions.
\Cref{sec:usage} gives usage recommendations and best practices.
\Cref{sec:discussion} concludes.

%% ============================================================================
\section{Theoretical Background}
\label{sec:theory}

\subsection{Mutual Information}
\label{sec:theory:mi}

The mutual information between two random variables $X$ and $Y$ with joint
density $p(x,y)$ and marginal densities $p(x)$, $p(y)$ is

\begin{equation}
  \MI{X;Y}
  = \int p(x,y)\log\frac{p(x,y)}{p(x)p(y)}\,dx\,dy
  = \KL{p(x,y)}{p(x)p(y)}.
  \label{eq:mi}
\end{equation}

For continuous, high-dimensional distributions, the densities in
\cref{eq:mi} are unknown and impossible to estimate accurately by
direct density estimation.  \nmi{} instead reframes MI estimation as a
neural network optimisation problem.

\subsection{Neural MI Estimators}
\label{sec:theory:estimators}

The central idea of variational MI estimation is to introduce a learnable
\emph{critic function} $f_\theta(x,y)$ and use it to compute a tractable
lower bound on $\MI{X;Y}$.  Given a batch of $N$ paired samples
$\{(x_i, y_i)\}_{i=1}^N$, the critic is trained by maximising the
estimate.

\subsubsection{InfoNCE}
\label{sec:theory:infonce}

The InfoNCE (Noise-Contrastive Estimation) estimator
\citep{oord2018representation} is the workhorse of \nmi{}:

\begin{equation}
  \hat{I}_{\mathrm{NCE}}
  = \E{\frac{1}{N}\sum_{i=1}^N
      \left(f(x_i,y_i) - \log\sum_{j=1}^N e^{f(x_i,y_j)}\right)} + \log N.
  \label{eq:infonce}
\end{equation}

\textbf{Intuition.}  For each anchor $x_i$, the critic assigns a score
to every $y_j$ in the batch.  The objective maximises the score of the
true partner $y_i$ relative to $N-1$ ``negative'' samples, i.e.\ it solves
a $N$-way classification problem.

\textbf{Properties.}
\begin{itemize}[leftmargin=*]
  \item \emph{Low variance.}  InfoNCE is highly stable across random seeds.
  \item \emph{Bounded bias.}  The bound is provably limited by $\log N$ (the
        batch size), so InfoNCE cannot report $\hat{I} > \log N$.  In practice
        this is rarely a limitation.  \nmi{} warns the user when the estimate
        approaches $0.85\,\log N_{\mathrm{eval}}$ and suggests increasing
        \api{max\_eval\_samples} or switching to SMILE.
  \item \emph{Default estimator.}  InfoNCE is the default in \nmi{} for all
        modes except \api{dimensionality}, where SMILE is recommended.
\end{itemize}

\subsubsection{SMILE}
\label{sec:theory:smile}

The SMILE (Smoothed Mutual Information Lower-bound Estimator) of
\citet{song2020understanding} reduces bias by clipping the normalising term:

\begin{equation}
  \hat{I}_{\mathrm{SMILE}}
  = \E{f(x,y)}
    - \log\E{e^{\mathrm{clip}(f(x,y'),\,\tau)}},
  \label{eq:smile}
\end{equation}

where $\mathrm{clip}(\cdot,\tau)$ clamps values to $[-\tau,\tau]$.

\textbf{Properties.}
\begin{itemize}[leftmargin=*]
  \item \emph{Lower bias.}  Clipping prevents a few ``easy'' negatives from
        dominating the log-normaliser, yielding estimates much closer to the
        true MI when it is large.  SMILE is not bounded by $\log N$.
  \item \emph{Moderate variance.}  The reduction in bias comes at the cost of
        slightly higher variance.  A clip value of $\tau=5$ is a robust default.
\end{itemize}

\subsubsection{Additional Estimators}
\label{sec:theory:other}

\nmi{} also implements the following estimators, accessible via
\api{estimator\_name} in \api{base\_params}:

\begin{itemize}[leftmargin=*]
  \item \textbf{TUBA} \citep{poole2019variational}: Tractable Unnormalized
        Barber-Agakov bound.
  \item \textbf{NWJ} \citep{poole2019variational}: Nguyen-Wainwright-Jordan
        bound; equivalent to TUBA shifted by 1.
  \item \textbf{JS-fGAN} \citep{poole2019variational}: Jensen-Shannon
        $f$-GAN lower bound based on the softplus activation.
\end{itemize}

\subsection{Variational Embeddings}
\label{sec:theory:variational}

Standard embedding networks map each input $x$ to a deterministic vector
$z = g_\theta(x)$.  \nmi{} additionally supports a variational approach in
which the encoder parameterises a Gaussian posterior:

\begin{equation}
  (\mu_x, \sigma_x) = g_\theta(x),
  \qquad
  z_x \sim \mathcal{N}(\mu_x, \sigma_x^2\,\mathbf{I}).
  \label{eq:varembed}
\end{equation}

The training loss combines the KL divergence from the posterior to a
standard-normal prior with the MI objective:

\begin{equation}
  \mathcal{L}
  = \KL{q(z_x|x)}{\mathcal{N}(0,\mathbf{I})}
  + \KL{q(z_y|y)}{\mathcal{N}(0,\mathbf{I})}
  - \beta\,\hat{I}_{\mathrm{estimator}}(Z_X;\,Z_Y).
  \label{eq:varloss}
\end{parameter}
\label{eq:varloss}
\end{equation}

The default $\beta = 1024$ makes MI maximisation strongly dominate over the
KL prior, which acts only as a mild regulariser on the embedding
distributions.  Decreasing $\beta$ increases the influence of the prior;
$\beta \ll 1$ can collapse embeddings toward the prior and reduce estimated MI.

\subsection{Finite-Sample Bias}
\label{sec:theory:bias}

Any MI estimate from a finite dataset of $N$ samples suffers from a
systematic upward bias because the model exploits random noise in the data.
Theoretically, for large $N$:

\begin{equation}
  I_{\mathrm{est}}(N) \approx I_{\mathrm{true}} + \frac{a}{N} + O(N^{-2}).
  \label{eq:bias}
\end{equation}

This linear relationship between the expected bias and $1/N$ is the key
insight exploited by the \api{rigorous} mode.

\subsection{Rigorous Bias Correction via Subsampling}
\label{sec:theory:rigorous}

The \api{rigorous} mode automates a three-step bias-correction procedure based
on \cref{eq:bias}:

\begin{enumerate}[leftmargin=*]
  \item \textbf{Subsampling.}  For each value of $\gamma \in \{1,2,\ldots,10\}$,
        the $N$-sample dataset is divided into $\gamma$ equal chunks.  Models
        are trained independently on each chunk of size $N/\gamma$.

  \item \textbf{Linear regression.}  Mean MI estimates at each $\gamma$ are
        plotted against $1/\gamma \propto 1/N$.  A weighted least-squares
        (WLS) regression is fitted in $1/\gamma$ space.

  \item \textbf{Extrapolation.}  The y-intercept of the fit—corresponding to
        $\gamma\to 0$ (i.e.\ $N\to\infty$)—is the bias-corrected MI estimate.
        The confidence interval of the intercept (default: 68\%, corresponding
        to $\pm 1\sigma$) provides the error bars.
\end{enumerate}

\textbf{Quadratic curvature filter.}  At very large $\gamma$ (very small
subsets), finite-sample effects and network under-fitting introduce curvature
that violates the linearity assumption.  \nmi{} automatically detects this by
fitting a quadratic polynomial and testing whether the ratio of the quadratic
to linear coefficient exceeds a configurable threshold
$\delta_{\mathrm{threshold}}$ (default 0.10).  Any $\gamma$ point that fails
this test is discarded iteratively from the largest $\gamma$ downward.  A
minimum of $\gamma_{\min}=5$ points must survive; if fewer remain the result
is flagged as \emph{unreliable}.

\textbf{Master-permutation subsampling.}  To ensure that the regression fits
pure $N$-dependent bias rather than sampling-noise variation, all gamma-levels
for a given hyperparameter combination share a \emph{single master
permutation}: the $\gamma=2$ subsets are literally the two halves of the
$\gamma=1$ dataset.

\subsection{Latent Dimensionality via Spectral Analysis}
\label{sec:theory:dim}

The \emph{intrinsic dimensionality} of the shared information between two
datasets $X$ and $Y$ can be measured without sweeping the bottleneck dimension,
by training an over-parameterised Hybrid Critic (see \cref{sec:arch:critics})
and analysing the spectrum of the cross-covariance of the learned embeddings.

\subsubsection{The Saturation Hypothesis}

When the bottleneck dimension $k_z$ is safely larger than the true
dimensionality $d^*$, the network concentrates all shared information into
a compact $d^*$-dimensional subspace of the latent space rather than
distributing it diffusely across all $k_z$ dimensions.

\subsubsection{Embedding Whitening}

Before computing the cross-covariance, \nmi{} whitens the embeddings.  The
default \api{spectral\_whitening='std'} standardises each embedding
dimension by its empirical standard deviation:

\begin{equation}
  \tilde{Z}_{X,i} = \frac{Z_{X,i}}{\mathrm{std}(Z_{X,i})},
  \qquad
  \tilde{Z}_{Y,i} = \frac{Z_{Y,i}}{\mathrm{std}(Z_{Y,i})}.
  \label{eq:whitening}
\end{equation}

This prevents dimensions with accidentally large variance from dominating
the SVD spectrum.  An alternative ZCA whitening (\api{'zca'}) is also
supported.

\subsubsection{Cross-Covariance and Participation Ratio}

Given $N$ test-set embeddings $Z_X \in \R^{N\times k_z}$ and
$Z_Y \in \R^{N\times k_z}$, the centred cross-covariance matrix is

\begin{equation}
  C_{XY} = \frac{1}{N-1}(Z_X - \bar{Z}_X)^\top(Z_Y - \bar{Z}_Y)
  \in \R^{k_z\times k_z}.
  \label{eq:crosscov}
\end{equation}

The singular value decomposition $C_{XY} = U\Sigma V^\top$ yields singular
values $\{\sigma_i\}_{i=1}^{k_z}$.  \nmi{} reports two variants of the
Participation Ratio (PR) \citep{gulati2026mutual}:

\begin{align}
  \mathrm{PR}_{\mathrm{sv}}
  &= \frac{\left(\sum_i \sigma_i\right)^2}{\sum_i \sigma_i^2},
  \label{eq:pr_sv}\\[4pt]
  \mathrm{PR}_{\mathrm{cov}}
  &= \frac{\left(\sum_i \sigma_i^2\right)^2}{\sum_i \sigma_i^4}.
  \label{eq:pr_cov}
\end{align}

Both are continuous measures of the effective number of shared dimensions.
$\mathrm{PR}_{\mathrm{cov}}$ weights each dimension by its eigenvalue
$\lambda_i = \sigma_i^2$, amplifying the contrast between large and small
singular values; it is therefore more sensitive to the true rank of the
representation.  In general, $\mathrm{PR}_{\mathrm{cov}} \le
\mathrm{PR}_{\mathrm{sv}}$.

\subsubsection{Interaction vs.\ Intrinsic Dimensionality}

\begin{itemize}[leftmargin=*]
  \item \textbf{Interaction dimensionality} ($y\_data$ provided): the PR of
        $C_{XY}$ directly yields the number of shared dimensions between two
        distinct populations.
  \item \textbf{Intrinsic dimensionality} (only $x\_data$): \nmi{} splits
        $x\_data$ channels into two non-overlapping halves and computes the
        interaction dimensionality between them, repeating over multiple random
        splits (default $n\_splits=5$) and averaging the PR.
\end{itemize}

\subsection{Spike-Timing Precision}
\label{sec:theory:precision}

The temporal precision of a neural code can be measured by determining how
much information is lost when spike times are progressively corrupted.
\nmi{} implements a \emph{Train Once, Evaluate Many} paradigm:

\begin{enumerate}[leftmargin=*]
  \item Train a baseline critic on uncorrupted data to convergence.
  \item Freeze all model weights.
  \item For each precision level $\tau$ in a user-specified grid, corrupt the
        target variable (\api{corrupt\_target}: \api{'x'}, \api{'y'}, or
        \api{'both'}) by one of two methods:
        \begin{itemize}
          \item \textbf{Deterministic rounding} (default):
                $\tilde{X} = \tau\lfloor X/\tau\rceil$, where $\lfloor\cdot\rceil$
                denotes rounding to the nearest integer.  Requires a single
                forward pass per $\tau$.
          \item \textbf{Additive uniform noise}:
                $\tilde{X} = X + \varepsilon$,
                $\varepsilon \sim U(-\tau/2,\,\tau/2)$.
                Averaged over $n\_noise\_samples=50$ realisations per $\tau$ for
                stability.
        \end{itemize}
  \item Evaluate the frozen critic on the corrupted data at each $\tau$,
        obtaining a degradation curve $\hat{I}(\tilde{X}^\tau;\,Y)$.
\end{enumerate}

The \emph{precision threshold} $\tau^*$ is the smallest $\tau$ at which

\begin{equation}
  \hat{I}(\tilde{X}^{\tau^*};\,Y) < \rho\cdot\hat{I}(X;\,Y),
  \label{eq:taustar}
\end{equation}

where $\rho$ is the \api{threshold\_ratio} (default 0.9).  Multiple ratios
can be specified simultaneously to characterise the full degradation profile.

\subsection{Transfer Entropy}
\label{sec:theory:te}

Transfer entropy (TE) from $X$ to $Y$ is defined via the conditional MI
chain rule \citep{schreiber2000measuring}:

\begin{equation}
  \mathrm{TE}(X\to Y)
  = I(x_{\mathrm{past}}, y_{\mathrm{past}};\, y_{\mathrm{future}})
    - I(y_{\mathrm{past}};\, y_{\mathrm{future}}).
  \label{eq:te}
\end{equation}

Past and future windows are constructed from the raw time series using
sliding windows of length \api{history\_window} and
\api{prediction\_horizon}, respectively.  Both MI terms in \cref{eq:te}
are estimated independently with \api{ParameterSweep}.

\subsection{Conditional Mutual Information}
\label{sec:theory:cmi}

The conditional MI between $X$ and $Y$ given a conditioning variable $Z$
is computed via the chain rule:

\begin{equation}
  I(X;\,Y\,|\,Z)
  = I(X,Z;\,Y) - I(Z;\,Y).
  \label{eq:cmi}
\end{equation}

$Z$ is concatenated with $X$ along the channel dimension before estimation,
so both terms in \cref{eq:cmi} use the same critic architecture and training
hyperparameters.

%% ============================================================================
\section{Library Architecture}
\label{sec:architecture}

\subsection{Package Overview}

\nmi{} is structured as a Python package with the following top-level
modules:

\begin{center}
\begin{tabular}{lp{10cm}}
\toprule
\textbf{Module / Subpackage} & \textbf{Responsibility} \\
\midrule
\api{neural\_mi.run} & Public entry point: \api{nmi.run()}, \api{nmi.extract\_embeddings()} \\
\api{neural\_mi.results} & \api{Results} container, \api{summary()}, \api{plot()}, persistence \\
\api{neural\_mi.defaults} & Schema validation for \api{base\_params} and mode kwargs \\
\api{neural\_mi.models} & Embedding networks and critic architectures \\
\api{neural\_mi.estimators} & Differentiable MI lower bounds (InfoNCE, SMILE, etc.) \\
\api{neural\_mi.training} & \api{Trainer} class: training loop, early stopping, spectral tracking \\
\api{neural\_mi.data} & Dataset classes and processor pipeline \\
\api{neural\_mi.analysis} & One submodule per analysis mode \\
\api{neural\_mi.generators} & Synthetic data generators \\
\api{neural\_mi.visualize} & Plotting utilities (bias-correction plots, etc.) \\
\api{neural\_mi.embeddings\_io} & Embedding extraction and model serialisation \\
\bottomrule
\end{tabular}
\end{center}

\subsection{Embedding Networks}
\label{sec:arch:embeddings}

All embedding networks inherit from \api{BaseEmbedding} and map an input
tensor of shape $(B,C,W)$ to a vector of shape $(B,d_z)$, where $B$ is
the batch size, $C$ the number of channels, $W$ the window length, and
$d_z$ the embedding dimension.

\subsubsection{MLP}
A standard multi-layer perceptron that flattens the input and passes it
through $n\_layers$ hidden blocks.  Each block consists of a
spectrally-normalised linear layer, optional normalisation
(\api{norm\_layer} $\in \{$\api{'batch'}, \api{'layer'}, \api{None}$\}$),
a configurable activation
(\api{'relu'}, \api{'tanh'}, \api{'leaky\_relu'}, \api{'silu'}),
and optional dropout.  Spectral normalisation is applied by default
(\api{use\_spectral\_norm=True}) to the hidden layers only; the output
linear layer is left unnormalised to preserve the full range of the
embedding.  \textbf{Recommended default for static or pre-windowed data.}

\subsubsection{CNN1D}
A stack of 1-D convolutional layers with odd kernel sizes (enforcing
``same'' padding to preserve sequence length), followed by
\api{AdaptiveAvgPool1d(1)} global average pooling and a two-layer MLP
projection.  Using \api{AdaptiveAvgPool1d} ensures the model is fully
picklable at construction time (important for multiprocessing workers).
\textbf{Recommended for time-series data where local temporal patterns
matter.}

\subsubsection{GRU and LSTM}
Recurrent networks that process the input as a sequence of length $W$
over $C$ feature channels.  Both support \api{bidirectional=True}.  The
embedding is extracted from the hidden state of the last layer (last
forward hidden state for unidirectional; concatenation of last forward
and backward hidden states for bidirectional).

\subsubsection{TCN}
A Temporal Convolutional Network consisting of dilated residual blocks
(\api{TemporalBlock}).  The dilation factor doubles with each layer
($2^i$ for layer $i$), giving an exponentially growing receptive field
of size $2^{n\_layers}(k-1)+1$, where $k$ is the kernel size.  The
embedding is taken from the last time step of the final layer.
\textbf{Recommended when long-range temporal context is important but
recurrent gradient flow is unstable.}

\subsubsection{Transformer}
A standard Transformer encoder with sinusoidal positional encoding.
The input is projected to \api{hidden\_dim} (\api{d\_model}), which must
be divisible by the number of attention heads \api{nhead}.  The embedding
is the output of the first token (CLS-style).  Xavier initialisation is
applied uniformly across all layers.

\subsubsection{VarMLP}
A variational counterpart to the standard MLP.  The network outputs both
a mean $\mu$ and a log-variance $\log\sigma^2$ vector.  During training
it applies the reparameterisation trick to sample $z = \mu + \sigma\varepsilon$
and returns the KL divergence alongside the embedding.  During evaluation
it returns $\mu$ and a zero KL loss.  This is the embedding used when
\api{use\_variational=True}.

\subsubsection{Shared Encoder}
When \api{shared\_encoder=True} (the default in \api{dimensionality} mode),
$X$ and $Y$ share a single embedding network.  This is
theoretically justified when both variables are split views of the same
data source, and reduces the parameter count by half.

\subsection{Critic Architectures}
\label{sec:arch:critics}

\nmi{} provides three critic types, controlled by the \api{critic\_type}
parameter.

\subsubsection{Separable Critic}

Embeds $X$ and $Y$ independently and computes the score as a dot product:
\begin{equation}
  f(x_i, y_j) = z_x(i)^\top z_y(j).
  \label{eq:sep}
\end{equation}
The score matrix $F \in \R^{N\times N}$ has $F_{ij} = z_x(i)^\top z_y(j)$,
computed as a single matrix multiplication $F = Z_X Z_Y^\top$.  This is
extremely efficient and is the \textbf{recommended default for standard MI
estimation}.  The dot-product geometry requires the true signal to live in
a well-aligned subspace of the embedding space; increasing
\api{embedding\_dim} helps when the signal is complex.

\subsubsection{Hybrid Critic}

Embeds $X$ and $Y$ independently to $\R^{d_z}$ each, concatenates the
pair into a $2d_z$-dimensional vector, and passes it through a two-layer
MLP decision head to produce a scalar score:
\begin{equation}
  f(x_i, y_j) = \mathrm{MLP}_\phi\bigl([z_x(i);z_y(j)]\bigr).
  \label{eq:hybrid}
\end{equation}
Because the decision head has unrestricted geometric flexibility, the
Hybrid Critic can detect \emph{any} nonlinear dependency that is captured
by the embeddings, regardless of their alignment.  This is
\textbf{required for the \api{dimensionality} mode} and is automatically
enforced there.

The full $N^2$ score matrix is computed in memory-safe chunks of size
\api{max\_n\_batches} to prevent out-of-memory errors on large batches.

\subsubsection{Concat Critic}

Concatenates the \emph{raw} inputs $x$ and $y$ before passing them through
a single network.  No separate embedding networks are used.  This is the
most flexible architecture but scales quadratically with input dimension
and is generally not recommended for high-dimensional data.

\subsection{Training}
\label{sec:arch:training}

The \api{Trainer} class encapsulates the complete training loop.

\subsubsection{Data Splitting}

Three modes are supported, with the following priority order:

\begin{enumerate}[leftmargin=*]
  \item \textbf{Custom indices} (\api{train\_indices}, \api{test\_indices}):
        user-provided arrays; all other split parameters are ignored.
  \item \textbf{Random split} (\api{split\_mode='random'}):
        a standard shuffled split of the $N$ indices, appropriate for IID data.
  \item \textbf{Blocked split} (\api{split\_mode='blocked'}, default):
        samples $k$ non-overlapping contiguous blocks from the full sequence to
        form the test set.  Block start positions are drawn using a minimum-distance
        algorithm to prevent consecutive blocks.  A configurable gap
        (\api{split\_gap\_fraction}) around each test block is excluded from the
        training set to reduce autocorrelation leakage.
\end{enumerate}

\subsubsection{Training Loop}

Each epoch performs the following steps:
\begin{enumerate}[leftmargin=*]
  \item \textbf{Stochastic gradient step.}  Training indices are randomly
        shuffled and split into batches.  The loss is
        $\mathcal{L} = -\hat{I}_{\mathrm{estimator}}$ (or
        $\mathcal{L} = \mathrm{KL} - \beta\hat{I}$ for variational models).
        Optional gradient clipping (\api{gradient\_clip\_val}) is applied
        before the optimiser step.

  \item \textbf{Validation evaluation.}  The model is evaluated on at most
        \api{max\_eval\_samples} held-out samples.  If the test set exceeds this
        limit, a single random subset of exactly \api{max\_eval\_samples} is
        drawn (unbiased, reduced-variance estimate).

  \item \textbf{Smoothing.}  The validation MI history is smoothed by a
        median filter (\api{median\_window}) followed by a Gaussian filter
        (\api{smoothing\_sigma}).  A custom smoothing function can be supplied.

  \item \textbf{Early stopping.}  If the smoothed MI does not improve by at
        least \api{min\_improvement} (relative) for \api{patience} consecutive
        epochs, training is halted.  The best model state is saved in memory
        via \api{copy.deepcopy} (no disk I/O).

  \item \textbf{Temporal augmentation.}  If the dataset is a
        \api{PairedTemporalDataset}, the window start positions are randomly
        shifted each epoch (linearly increasing from zero to
        \api{window\_size} over \api{epochs\_to\_max\_shift} epochs), so the
        model sees a wide range of temporal contexts.

  \item \textbf{Learning-rate scheduling.}  An optional PyTorch LR scheduler
        is stepped after each epoch.  Supported schedulers: \api{'cosine'},
        \api{'step'}, \api{'plateau'}, \api{'cosine\_warmup'} (linear warmup
        followed by cosine annealing), or any custom
        \api{torch.optim.lr\_scheduler} instance.
\end{enumerate}

\textbf{NaN handling.}  If the MI evaluation returns NaN, the epoch is
skipped for the purpose of early stopping and a warning is logged.  Three
consecutive NaN values raise a \api{TrainingError} to prevent silent
numerical instability.

\textbf{InfoNCE ceiling detection.}  After training, if the final test MI
exceeds $0.85\,\log(N_{\mathrm{eval}})$, a warning is issued.

\subsubsection{Optimisers}

All standard PyTorch optimisers are accessible via the string names
\api{'adam'} (default), \api{'adamw'}, \api{'sgd'}, \api{'rmsprop'},
\api{'adagrad'}.  A custom optimiser class can also be passed directly.
The learning rate and additional optimiser parameters are set via
\api{learning\_rate} and \api{optimizer\_params}.

\subsubsection{Spectral Metric Tracking}

When \api{track\_spectral\_metrics=True}, the trainer computes the
cross-covariance spectrum of the embeddings at every epoch and stores the
history in the returned dictionary.  At the end of training, the final
Participation Ratio is always computed and returned.  Setting
\api{spectral\_output='all'} returns all spectral metrics
(PR-covariance, PR-singular, effective rank, spectral entropy);
the default \api{'default'} returns only $\mathrm{PR}_{\mathrm{sv}}$ as
\api{participation\_ratio}.  Setting \api{return\_spectrum=True}
additionally stores the raw singular-value array.

%% ============================================================================
\section{Data Handling}
\label{sec:data}

\subsection{Processor Classes}
\label{sec:data:processors}

\nmi{} ships three processor classes that convert raw recordings to the
canonical tensor format $(N_{\mathrm{samples}},\,C,\,W)$, where
$N_{\mathrm{samples}}$ is the number of windows, $C$ is the number of
channels, and $W$ is the window length.

\subsubsection{ContinuousProcessor}
For continuously sampled signals such as LFP, EEG, calcium imaging, or
kinematics.  The processor divides the raw time series into non-overlapping
windows of length \api{window\_size} samples.  The \api{time\_shift}
parameter controls an optional initial offset.  The sample rate
(\api{sample\_rate}) is used by the \api{lag} mode to express lag ranges
in seconds rather than samples.

\subsubsection{SpikeProcessor}
Converts raw spike-time lists (one list per neuron) to binned
\mbox{$(0/1)$} or count arrays inside windows of length
\api{window\_size} (in seconds).  Additional parameters include
\api{max\_spikes\_per\_window} and \api{n\_seconds}.  The resulting
tensor has shape $(N_{\mathrm{windows}},\,N_{\mathrm{neurons}},\,W)$.

\subsubsection{CategoricalProcessor}
Encodes discrete behavioural or stimulus state sequences.  Each category
is mapped to a one-hot vector or an integer index, depending on the
\api{encoding} parameter.  The output has shape
$(N,\,C,\,N_{\mathrm{categories}})$.

\subsubsection{Fast Path}
When both \api{processor\_type\_x} and \api{processor\_type\_y} are
\api{None}, \nmi{} bypasses the processor pipeline entirely and converts
the user-supplied arrays directly to PyTorch tensors.  This fast path
avoids one \api{PairedDataset} allocation per call and is the appropriate
choice when data are already pre-processed.

\subsection{Dataset Classes}
\label{sec:data:datasets}

\begin{itemize}[leftmargin=*]
  \item \textbf{PairedDataset}: wraps a pair of static tensors $(X, Y)$.
        Stores both a mutable working copy (\api{data}) and an immutable
        master copy (\api{data\_master}).  Noise and corruption operations
        derive their mask from the master copy to prevent compound effects
        across repeated applications.

  \item \textbf{PairedTemporalDataset}: extends \api{PairedDataset} with
        a \api{WindowManager} that maintains the set of valid window start
        positions after any time shift.  Used for continuous and spike
        data.

  \item \textbf{SubsetView}: a lightweight indexing wrapper that makes an
        arbitrary subset of a dataset appear as a standalone dataset,
        without copying data.  Used to implement train/test splits without
        allocating new tensors.

  \item \textbf{ContinuousWindowDataset}, \textbf{SpikeWindowDataset},
        \textbf{CategoricalWindowDataset}: lower-level dataset classes that
        store the raw recordings and perform on-the-fly windowing plus
        optional noise augmentation.
\end{itemize}

\subsection{Multi-modal Alignment}
\label{sec:data:multimodal}

When $X$ and $Y$ have different sampling rates, recording durations, or
window sizes, \nmi{} aligns them by truncating to the shorter stream.
Different window sizes for $X$ and $Y$ are not supported within a single
run (the \api{window\_size} parameter in \api{processor\_params\_y} is
ignored with a warning if it differs from \api{processor\_params\_x}).

%% ============================================================================
\section{Analysis Modes}
\label{sec:modes}

All modes are accessed via the unified entry point:
\begin{lstlisting}
import neural_mi as nmi
result = nmi.run(
    x_data, y_data,
    mode='rigorous',       # one of 9 modes
    base_params={...},     # model architecture + training schedule
    **mode_specific_kwargs
)
\end{lstlisting}

The function returns a \api{Results} object (see \cref{sec:results}).  The
complete parameter reference for \api{nmi.run()} is given in
\cref{app:run_params}.

\subsection{Estimate (\texttt{mode='estimate'})}
\label{sec:modes:estimate}

\textbf{Purpose.}  Quick point estimate of $\MI{X;Y}$ for initial
exploration.

\textbf{Method.}  Runs a single training task via the \api{ParameterSweep}
infrastructure with an empty sweep grid.

\textbf{Key parameters.}
\begin{itemize}[leftmargin=*]
  \item All \api{base\_params}: architecture, training schedule, estimator.
  \item \api{permutation\_test} (bool): whether to compute a null
        distribution by randomly permuting $Y$ \api{n\_permutations} times
        (default 100).  The permutation test is available for all modes
        except \api{rigorous} and \api{precision}.
  \item \api{output\_units}: \api{'nats'} (default) or \api{'bits'}.
\end{itemize}

\textbf{Returns.}  A \api{Results} object with \api{mi\_estimate} set to
\api{train\_mi} from the single training run.

\subsection{Sweep (\texttt{mode='sweep'})}
\label{sec:modes:sweep}

\textbf{Purpose.}  Parallelised sweep over any combination of
hyperparameters.  Used to explore model capacity, learning rates, random
seeds, window sizes, or data-processing parameters.

\textbf{Method.}  The \api{ParameterSweep} class takes a \api{base\_params}
dictionary and a \api{sweep\_grid} dictionary (mapping parameter names to
lists of values).  It generates the full Cartesian product of values,
runs each combination as an independent training task, and aggregates the
results into a \api{DataFrame}.

\textbf{Key parameters.}
\begin{itemize}[leftmargin=*]
  \item \api{sweep\_grid} (dict): e.g.\ \api{\{'run\_id': range(5),
        'embedding\_dim': [16, 32, 64]\}}.
  \item \api{n\_workers} (int): number of parallel processes (default 1).
  \item \api{processor\_type\_x/y}: if either is set to a
        \emph{processor name} (rather than a class), the sweep will
        include a grid over processor parameters.
\end{itemize}

\textbf{Returns.}  \api{Results} with \api{mi\_estimate} set to the mean
of all \api{train\_mi} values across the sweep, and a \api{DataFrame} in
\api{details['dataframe']}.

\subsection{Lag (\texttt{mode='lag'})}
\label{sec:modes:lag}

\textbf{Purpose.}  Nonlinear cross-correlation: find the temporal offset
$\ell^* = \arg\max_\ell \hat{I}(X;\,Y_{\ell})$, where $Y_\ell$ is $Y$
shifted by $\ell$ samples (or seconds if \api{sample\_rate} is known).

\textbf{Method.}  For each lag in \api{lag\_range}:
\begin{enumerate}[leftmargin=*]
  \item Shift $Y$ by $\ell$ samples (truncating both $X$ and $Y$ to the
        overlapping segment).
  \item Run a full, independent MI estimation on the shifted pair.
\end{enumerate}
All lag tasks are dispatched to \api{n\_workers} parallel processes.

Optional \api{equalize\_n=True} truncates all lags to the same number of
samples as the largest lag (preventing the apparent MI rise from larger
sample sizes at smaller lags).

\textbf{Key parameters.}
\begin{itemize}[leftmargin=*]
  \item \api{lag\_range}: a Python \api{range} or list of integers
        (samples) or floats (seconds if \api{sample\_rate} is provided).
  \item \api{sample\_rate} (float, optional): if provided (in
        \api{processor\_params\_x}), lags are interpreted in seconds.
  \item \api{equalize\_n} (bool): equalise sample count across lags
        (default \api{False}).
\end{itemize}

\textbf{Returns.}  \api{Results} with \api{mi\_estimate} set to the best
MI across lags, and a \api{DataFrame} in \api{details['dataframe']}
containing columns \api{lag}, \api{train\_mi}, \api{n\_windows} per lag.

\subsection{Precision (\texttt{mode='precision'})}
\label{sec:modes:precision}

\textbf{Purpose.}  Quantify the temporal resolution at which a neural
code carries information, yielding an empirical precision threshold
$\tau^*$ in the same units as the data.

\textbf{Method.}  See \cref{sec:theory:precision}.

\textbf{Key parameters.}
\begin{itemize}[leftmargin=*]
  \item \api{tau\_grid} (list of float): the precision levels to test.
  \item \api{corrupt\_target}: \api{'x'} (default), \api{'y'}, or
        \api{'both'}.
  \item \api{corruption\_method}: \api{'rounding'} (default) or
        \api{'noise'}.
  \item \api{n\_noise\_samples} (int, default 50): number of stochastic
        noise realisations per $\tau$ when \api{corruption\_method='noise'}.
  \item \api{threshold\_ratio} (float or list): the fraction of baseline MI
        that defines $\tau^*$.  A list produces multiple thresholds
        simultaneously (primary threshold is the first value).
\end{itemize}

\textbf{Returns.}  \api{Results} with:
\begin{itemize}[leftmargin=*]
  \item \api{mi\_estimate}: baseline MI at zero corruption.
  \item \api{details['precision\_tau']}: the primary threshold $\tau^*$.
  \item \api{details['baseline\_mi']}: MI at zero corruption (same as
        \api{mi\_estimate}).
  \item \api{details['precision\_thresholds']}: full dict mapping each
        ratio to its $\{\api{precision\_tau},\,\api{threshold\_value}\}$.
  \item \api{details['dataframe']}: DataFrame with columns \api{tau},
        \api{train\_mi}, \api{train\_mi\_std}.
\end{itemize}

\subsection{Rigorous (\texttt{mode='rigorous'})}
\label{sec:modes:rigorous}

\textbf{Purpose.}  Debiased MI estimate with a confidence interval,
correcting for the $O(1/N)$ finite-sample bias.

\textbf{Method.}  See \cref{sec:theory:rigorous}.

\textbf{Key parameters.}
\begin{itemize}[leftmargin=*]
  \item \api{gamma\_range} (range, default \api{range(1,11)}): the data
        fraction denominators.
  \item \api{delta\_threshold} (float, default 0.10): quadratic curvature
        threshold for linearity check.
  \item \api{min\_gamma\_points} (int, default 5): minimum surviving
        gamma points for the estimate to be flagged as reliable.
  \item \api{confidence\_level} (float, default 0.68): confidence level
        for the error bars (0.68 $\approx 1\sigma$).
  \item \api{min\_reliable\_samples} (int, default 1000): if the smallest
        data chunk is below this size, a warning is issued.
  \item \api{n\_workers}: parallelisation over gamma $\times$ sweep-grid tasks.
\end{itemize}

\textbf{Returns.}  \api{Results} with:
\begin{itemize}[leftmargin=*]
  \item \api{mi\_estimate}: bias-corrected MI (intercept of WLS fit).
  \item \api{details['mi\_error']}: half-width of the confidence interval.
  \item \api{details['slope']}: slope of the linear fit (proportional to
        the bias rate $a$ in \cref{eq:bias}).
  \item \api{details['is\_reliable']}: boolean.
  \item \api{details['gammas\_used']}: list of $\gamma$ values that
        survived the linearity filter.
  \item \api{details['raw\_results\_df']}: DataFrame of all per-task results.
\end{itemize}

\subsection{Dimensionality (\texttt{mode='dimensionality'})}
\label{sec:modes:dimensionality}

\textbf{Purpose.}  Estimate the effective number of dimensions of the
shared information space between two populations (Interaction Dimensionality)
or of a single population's intrinsic structure (Intrinsic Dimensionality).

\textbf{Method.}  See \cref{sec:theory:dim}.

\textbf{Key parameters (mode-specific).}
\begin{itemize}[leftmargin=*]
  \item \api{split\_method}: \api{'random'} (default), \api{'spatial'},
        or \api{'temporal'}.  Only used for Intrinsic Dimensionality
        (single-dataset mode).
  \item \api{n\_splits} (int, default 5): number of independent splits
        (Intrinsic) or independent model fits (Interaction) to average over.
  \item \api{spectral\_mode}: \api{'summary'} (default, returns only
        $\mathrm{PR}_{\mathrm{sv}}$) or \api{'full'} (all spectral metrics
        plus the raw singular-value array).
  \item \api{embedding\_dim}: bottleneck dimension; auto-defaults to 64 if
        not specified.
  \item \api{shared\_encoder}: whether $X$ and $Y$ share embedding weights
        (auto-defaults to \api{True} in this mode).
\end{itemize}

\textbf{Forced overrides.}  \nmi{} always sets \api{critic\_type='hybrid'}
in this mode; user values are silently replaced with a logged info message.

\textbf{Returns.}  \api{Results} with:
\begin{itemize}[leftmargin=*]
  \item \api{mi\_estimate}: mean MI across splits.
  \item \api{details['dataframe']}: one row per split with columns
        \api{split\_id}, \api{train\_mi}, \api{participation\_ratio}.
  \item \api{details['participation\_ratio']}: mean PR (default: $\mathrm{PR}_{\mathrm{sv}}$).
  \item \api{details['participation\_ratio\_std']}: standard deviation across splits.
  \item When \api{spectral\_mode='full'}: additional columns
        \api{pr\_covariance}, \api{effective\_rank}, \api{spectral\_entropy},
        and the raw \api{spectrum} (singular values).
\end{itemize}

\subsection{Conditional (\texttt{mode='conditional'})}
\label{sec:modes:conditional}

\textbf{Purpose.}  Measure the direct relationship between $X$ and $Y$
after removing the effect of a known confound $Z$:
$I(X;\,Y\,|\,Z)$.

\textbf{Method.}  See \cref{sec:theory:cmi}.  Two independent
\api{ParameterSweep} runs estimate $I(XZ;\,Y)$ and $I(Z;\,Y)$; the
conditional MI is their difference.

\textbf{Key parameters.}
\begin{itemize}[leftmargin=*]
  \item \api{z\_data}: the conditioning variable tensor, shape
        $(N,\,C_Z,\,W)$.
  \item \api{z\_processor\_type} / \api{z\_processor\_params} /
        \api{z\_time}: optional processor for $Z$ (follows the same
        pipeline as $X$ and $Y$).
  \item All sweep-grid parameters are forwarded to both inner sweeps.
\end{itemize}

\textbf{Returns.}  \api{Results} with:
\begin{itemize}[leftmargin=*]
  \item \api{mi\_estimate}: $I(X;\,Y\,|\,Z)$.
  \item \api{details['mi\_xz\_y']}: $I(XZ;\,Y)$.
  \item \api{details['mi\_z\_y']}: $I(Z;\,Y)$.
  \item \api{details['raw\_xz\_y']}, \api{details['raw\_z\_y']}:
        raw sweep results for each component.
\end{itemize}

\subsection{Transfer (\texttt{mode='transfer'})}
\label{sec:modes:transfer}

\textbf{Purpose.}  Quantify the directed information flow from $X$ to $Y$
(and optionally from $Y$ to $X$) using transfer entropy.

\textbf{Method.}  See \cref{sec:theory:te}.  Sliding-window past/future
arrays are built internally from the raw 2-D time series.

\textbf{Key parameters.}
\begin{itemize}[leftmargin=*]
  \item \api{history\_window} (int): past window length in samples.
  \item \api{prediction\_horizon} (int, default 1): how many steps ahead
        to predict.
  \item \api{bidirectional\_te} (bool, default \api{False}): if
        \api{True}, also computes $\mathrm{TE}(Y\to X)$ and returns a
        directionality index.
  \item \api{x\_data} and \api{y\_data} must be 2-D tensors of shape
        $(T,\,C)$.
\end{itemize}

\textbf{Returns.}  \api{Results} with:
\begin{itemize}[leftmargin=*]
  \item \api{mi\_estimate}: $\mathrm{TE}(X\to Y)$.
  \item \api{details['te\_xy']}: same as \api{mi\_estimate}.
  \item \api{details['i\_xypast\_yfuture']}: $I(x_{\mathrm{past}},\,y_{\mathrm{past}};\,y_{\mathrm{future}})$.
  \item \api{details['i\_ypast\_yfuture']}: $I(y_{\mathrm{past}};\,y_{\mathrm{future}})$.
  \item \api{details['n\_samples']}: number of valid sliding-window
        samples.
  \item When \api{bidirectional\_te=True}: \api{details['te\_yx']},
        \api{details['directionality\_index']} $\in [-1,+1]$ (positive
        = net $X\to Y$ flow).
\end{itemize}

\subsection{Pairwise (\texttt{mode='pairwise'})}
\label{sec:modes:pairwise}

\textbf{Purpose.}  Rapidly compute all-to-all functional connectivity
matrices.

\textbf{Method.}  For each channel pair $(i,j)$:
\begin{itemize}[leftmargin=*]
  \item \textbf{Self-pairwise} ($y\_data=\text{None}$): estimates MI
        between every unique pair $(i,j)$ with $i<j$ from $x\_data$.
        The result is a symmetric matrix with diagonal set to zero.
  \item \textbf{Cross-pairwise} ($y\_data$ provided): estimates MI between
        every channel $i$ of $x\_data$ and every channel $j$ of $y\_data$,
        producing a full $(C_X \times C_Y)$ matrix.
\end{itemize}

Mean and standard deviation of the MI across any sweep repetitions are
stored as \api{mi\_mean} and \api{mi\_std}.

\textbf{Key parameters.}
\begin{itemize}[leftmargin=*]
  \item \api{pairs} (list of tuples, optional): explicit list of
        $(i,j)$ pairs to estimate, rather than all pairs.
  \item \api{n\_workers}: parallelisation across channel pairs.
  \item \api{permutation\_test}: available; triggers a warning about
        computational cost ($N_{\mathrm{pairs}} \times n\_permutations$ full
        training runs).
\end{itemize}

\textbf{Returns.}  \api{Results} with:
\begin{itemize}[leftmargin=*]
  \item \api{details['mi\_matrix']}: np.ndarray of shape
        $(C_X,C_X)$ (self-pairwise) or $(C_X,C_Y)$ (cross-pairwise).
  \item \api{details['dataframe']}: DataFrame with columns \api{ch\_x},
        \api{ch\_y}, \api{mi\_mean}, \api{mi\_std}.
  \item \api{details['n\_channels']}: int or (int, int).
\end{itemize}

%% ============================================================================
\section{Embedding Extraction}
\label{sec:embeddings}

After training, the learned embedding networks can be extracted and applied
to new data via \api{nmi.extract\_embeddings()}:

\begin{lstlisting}
embeddings = nmi.extract_embeddings(
    result,          # Results object from a previous nmi.run() call
    x_new,           # new data to embed
    y_new,           # optional
)
# embeddings['z_x'].shape == (N, embedding_dim)
# embeddings['z_y'].shape == (N, embedding_dim)
\end{lstlisting}

Internally this calls \api{critic.get\_embeddings()}, which uses the same
chunked evaluation pipeline as during training to prevent out-of-memory
errors on large datasets.  For \api{ConcatCritic}, where there are no
separate embedding networks, the raw (flattened) inputs are returned.

The set of architecture parameters required to reconstruct a saved model
is exposed as the module-level constant
\api{neural\_mi.analysis.task.\_BUILD\_PARAMS\_KEYS} (includes
\api{critic\_type}, all embedding network parameters, \api{dropout},
\api{norm\_layer}, \api{use\_spectral\_norm}, etc.).

%% ============================================================================
\section{Synthetic Data Generators}
\label{sec:generators}

The \api{nmi.generators} submodule provides eight synthetic generators for
testing and validating MI estimators.  All generators return
\api{(x,\,y)} tuples and accept a \api{use\_torch} parameter.

\begin{description}[style=nextline, leftmargin=1.5cm]

\item[\api{generate\_correlated\_gaussians}($n$, $d$, $I$)]
  Two $d$-dimensional correlated multivariate Gaussians with analytically
  known MI $I$ bits.  The cross-correlation coefficient is derived from the
  Gaussian MI formula:
  $\rho = \sqrt{1 - e^{-2I\ln 2 / d}}$.
  Ground truth: $I_{\mathrm{true}} = I$ bits.

\item[\api{generate\_nonlinear\_from\_latent}($n$, $d_z$, $d_{\mathrm{obs}}$, $I$)]
  A shared $d_z$-dimensional latent variable is created with known MI $I$.
  Two $d_{\mathrm{obs}}$-dimensional observed variables are obtained by
  passing it through independent random MLP projections with Softplus
  nonlinearities.  Tests the ability of the estimator to detect nonlinear
  dependencies in high-dimensional data.

\item[\api{generate\_temporally\_convolved\_data}($n$, $\ell$, $\sigma$)]
  $Y$ is a lagged, noisy copy of $X$: $Y_t = X_{t-\ell} + \varepsilon_t$,
  $\varepsilon_t \sim \mathcal{N}(0,\sigma^2)$.
  $X$ is a smoothed Gaussian random walk.
  Shape: $(n,1)$.  Used for basic lag-analysis validation.

\item[\api{generate\_correlated\_spike\_trains}($N_{\mathrm{neurons}}$,\ldots)]
  Two populations of Poisson spike trains where population $Y$ fires with a
  Gaussian-jittered delay $\delta$ after population $X$.

\item[\api{generate\_correlated\_categorical\_series}($n$, $C$, $K$, $p$)]
  Two correlated Markov-chain categorical series of length $n$, $C$
  channels, $K$ categories, with $Y$ matching $X$ with probability $p$.

\item[\api{generate\_event\_related\_data}($n$, $\ell$, $n_{\mathrm{ev}}$,\ldots)]
  $X$ contains $n_{\mathrm{ev}}$ sharp unit impulses at random times.
  $Y$ is a sine-wave response that begins $\ell$ samples after each
  impulse.  Tests precision of event-triggered analysis.

\item[\api{generate\_xor\_data}($n$, $\sigma$)]
  Classic XOR dataset: $X = (X_1,X_2) \in \{0,1\}^2$ (uniform binary),
  $Y = \mathrm{XOR}(X_1,X_2) + \varepsilon$.  MI between either marginal
  and $Y$ is zero, but the joint MI is high.  Tests detection of
  \emph{synergistic} dependencies.

\item[\api{generate\_windowed\_dependency\_data}($T$, $C$, $\tau$, $W$, $\sigma_n$)]
  Multi-channel autocorrelated $X$ (mixture of fast and slow AR processes).
  $Y_t$ is a causal rolling mean of $X$ over a window $W$, plus noise.
  Isolates window-size effects from lag effects.  Useful for sweeping
  \api{window\_size}.

\end{description}

Additional generators for tutorial use:
\api{generate\_linear\_data}, \api{generate\_nonlinear\_data},
\api{generate\_history\_data}, \api{generate\_full\_data}.

%% ============================================================================
\section{The Results Object}
\label{sec:results}

All analysis modes return a \api{Results} instance with a common interface.

\subsection{Attributes}

\begin{center}
\begin{tabular}{lp{9cm}}
\toprule
\textbf{Attribute} & \textbf{Description} \\
\midrule
\api{mode} & Analysis mode string. \\
\api{mi\_estimate} & Primary scalar MI estimate (mode-specific meaning). \\
\api{params} & Copy of all \api{base\_params} used. \\
\api{details} & Mode-specific dict with full numerical output. \\
\api{critic} & The trained \api{nn.Module} (if retained). \\
\midrule
\multicolumn{2}{l}{\textit{Mode-specific content of} \api{details}} \\
\midrule
\textit{estimate/sweep} & \api{dataframe}, \api{sweep\_grid} \\
\textit{rigorous} & \api{mi\_error}, \api{slope}, \api{is\_reliable}, \api{gammas\_used}, \api{raw\_results\_df} \\
\textit{lag}      & \api{dataframe} (lag, train\_mi, n\_windows), \api{best\_lag} \\
\textit{precision} & \api{baseline\_mi}, \api{precision\_tau}, \api{threshold\_ratio},
                     \api{precision\_thresholds}, \api{dataframe} \\
\textit{dimensionality} & \api{dataframe}, \api{participation\_ratio},
                          \api{participation\_ratio\_std} \\
\textit{conditional} & \api{mi\_xz\_y}, \api{mi\_z\_y}, \api{raw\_xz\_y}, \api{raw\_z\_y} \\
\textit{transfer} & \api{te\_xy}, \api{i\_xypast\_yfuture}, \api{i\_ypast\_yfuture},
                    optionally \api{te\_yx}, \api{directionality\_index} \\
\textit{pairwise} & \api{mi\_matrix}, \api{dataframe} (ch\_x, ch\_y, mi\_mean, mi\_std) \\
\textit{any mode}  & \api{null\_distribution} (if \api{permutation\_test=True}) \\
\bottomrule
\end{tabular}
\end{center}

\subsection{Methods}

\begin{description}[style=nextline, leftmargin=1.5cm]
  \item[\api{summary()}]
    Prints a human-readable summary to stdout.  The output is tailored to
    each mode: \textit{precision} reports baseline MI, threshold, and ratio;
    \textit{pairwise} reports the matrix shape and mean MI; \textit{conditional}
    reports the three MI components; \textit{transfer} reports both TE
    directions and the directionality index.

  \item[\api{plot()}]
    For \textit{rigorous} mode, generates the bias-correction plot (MI vs.\
    $1/\gamma$ with the extrapolation line and confidence interval).  For
    \textit{sweep}/\textit{lag} modes, generates a bar or line plot of MI
    vs.\ sweep parameter.  For \textit{precision}, plots the degradation
    curve.

  \item[\api{save(path=None) \to str}]
    Serialises the entire \api{Results} object to a pickle file.
    If \api{path=None}, an auto-timestamped filename is generated in the
    current directory.  A non-colliding suffix is appended automatically
    if the file already exists.  Returns the full path used.

  \item[\api{load(path) \to Results}]
    Class method.  Deserialises a \api{Results} object from a pickle file
    created by \api{save()}.  Raises \api{TypeError} if the file does not
    contain a \api{Results} instance.

  \item[\api{to\_json(path=None) \to str}]
    Exports a human-readable JSON snapshot.  Large objects (arrays with more
    than 100 elements, DataFrames, raw result lists) are replaced by a
    short summary string rather than serialised in full.  Useful for logging
    experiment metadata.

\end{description}

%% ============================================================================
\section{Usage Recommendations and Best Practices}
\label{sec:usage}

\subsection{Choosing an Estimator}

\begin{itemize}[leftmargin=*]
  \item Use \textbf{InfoNCE} (default) for all standard estimation tasks.
        It is stable, fast, and its ceiling at $\log N$ is rarely a
        practical limitation.
  \item Switch to \textbf{SMILE} when the true MI is expected to be very
        high (e.g.\ near the $\log N$ InfoNCE ceiling), or when using
        \api{dimensionality} mode where the MI signal within the bottleneck
        can be large.
  \item Set \api{clip=5.0} for SMILE (default); decrease it toward 1.0 for
        more aggressive bias reduction at the cost of higher variance.
\end{itemize}

\subsection{Choosing a Critic}

\begin{itemize}[leftmargin=*]
  \item \textbf{SeparableCritic} (default): suitable for nearly all tasks
        when the embedding dimension is large enough ($\ge 32$).  Very fast
        due to the matrix-multiplication score computation.
  \item \textbf{HybridCritic}: required for \api{dimensionality} mode.
        Also recommended when the dependence structure is highly nonlinear
        and cannot be captured by a dot product in the embedding space.
  \item \textbf{ConcatCritic}: avoid for high-dimensional inputs; no
        independent embeddings are produced, so it cannot be used with
        \api{extract\_embeddings()}.
\end{itemize}

\subsection{Choosing an Embedding Model}

\begin{itemize}[leftmargin=*]
  \item \textbf{MLP}: for pre-windowed, static, or low-temporal-structure
        data.  Fastest.
  \item \textbf{CNN1D}: for time-series data where local features (e.g.\
        oscillation cycles) are informative.
  \item \textbf{GRU/LSTM}: when the temporal order of features within a
        window is important.  Use when the window is long enough to benefit
        from recurrent context.
  \item \textbf{TCN}: a faster, parallelisable alternative to RNNs with
        exponentially growing receptive field.
  \item \textbf{Transformer}: for long windows where attention-based
        long-range dependencies are expected.  Requires $d_{\mathrm{model}}$
        to be divisible by \api{nhead}.
\end{itemize}

\subsection{Data Splitting}

\begin{itemize}[leftmargin=*]
  \item Use \api{split\_mode='blocked'} (default) for all temporally
        autocorrelated data (LFP, EEG, spike trains, calcium imaging).
        This prevents the model from being tested on samples whose
        neighbours were used for training.
  \item Use \api{split\_mode='random'} only for genuinely IID data (e.g.\
        the result of extracting non-overlapping trials).
  \item Increase \api{n\_test\_blocks} (default 5) if the data have
        long-range autocorrelations (e.g.\ slow cortical oscillations at
        0.1 Hz sampled at 1000 Hz).
  \item Consider \api{split\_gap\_fraction} $> 0$ when strong
        autocorrelation persists at the boundary between training and
        validation blocks.
\end{itemize}

\subsection{Regularisation}

\begin{itemize}[leftmargin=*]
  \item For small datasets ($N < 500$ samples), add dropout
        (\api{dropout=0.1--0.3}), layer normalisation
        (\api{norm\_layer='layer'}), or weight decay
        (\api{optimizer='adamw'}, \api{optimizer\_params=\{'weight\_decay': 1e-4\}}).
  \item Spectral normalisation (\api{use\_spectral\_norm=True}, on by
        default) stabilises training and should rarely be disabled.
\end{itemize}

\subsection{Parallelisation}

\begin{itemize}[leftmargin=*]
  \item Set \api{n\_workers} to the number of available CPU cores for all
        modes except \api{precision} (which is inherently sequential).
  \item Multiprocessing uses the \api{'spawn'} context to avoid fork-related
        issues with CUDA tensors.  On macOS/Windows this adds a small
        startup overhead per worker.
  \item Move data to CPU before passing it to parallel modes
        (\api{x\_data.cpu()}) to avoid CUDA shared-memory complications.
\end{itemize}

\subsection{Rigorous Mode}

\begin{itemize}[leftmargin=*]
  \item Each gamma requires multiple training runs; the total number of runs
        is $|\gamma\_\mathrm{range}| \times |\text{sweep combinations}|$.
        For default settings (10 gamma values, 1 combination), this is 10
        full training tasks plus 10 subsets.
  \item If computation budget is limited, restrict \api{gamma\_range} to
        e.g.\ \api{range(1,7)} at the cost of a wider confidence interval.
  \item Verify \api{is\_reliable=True} before trusting the debiased
        estimate; if \api{False}, collect more data or reduce the maximum
        $\gamma$.
\end{itemize}

\subsection{Precision Mode}

\begin{itemize}[leftmargin=*]
  \item Express \api{tau\_grid} in the same units as the raw data (e.g.\
        milliseconds for spike times in ms).
  \item Use \api{corruption\_method='rounding'} (default) whenever
        possible—it is deterministic and requires only one forward pass per
        $\tau$.
  \item Increase \api{n\_noise\_samples} to $\ge 100$ for smooth noise
        curves if the degradation is expected to be monotone but noisy.
\end{itemize}

%% ============================================================================
\section{Discussion}
\label{sec:discussion}

\nmi{} provides a unified, scientifically rigorous framework for neural
mutual information estimation in neuroscience.  By combining flexible
neural network architectures with automated bias correction, principled
data splitting, and domain-specific data processors, it addresses the
main practical barriers to using MI in typical systems neuroscience
workflows.

\textbf{Relation to prior work.}  Existing neural MI implementations
\citep{oord2018representation,song2020understanding,belghazi2018mine}
provide point estimates for IID data with no bias correction, no
neuroscience-native data handling, and no unified interface.  \nmi{}
incorporates the $O(1/N)$ bias-correction framework of
\citet{abdelaleem2025accurate} and the spectral dimensionality approach
of \citet{gulati2026mutual} within a single maintained library.

\textbf{Limitations.}
\begin{itemize}[leftmargin=*]
  \item Neural MI estimators are lower bounds; the true MI can always exceed
        the estimate.  The $\log N$ InfoNCE ceiling can be a practical
        constraint for very large MI or very small batch sizes.
  \item The bias-correction extrapolation assumes the $O(1/N)$ regime
        applies at the sample sizes of interest; very small datasets may not
        enter the linear regime within the default $\gamma$ range.
  \item Transfer entropy as implemented here uses a chain-rule decomposition
        that requires two independent MI estimates; their individual
        finite-sample biases partially cancel, but the combination can
        still be noisy.
  \item All modes except \api{precision} train a fresh model per task and
        may be slow for large sweep grids without sufficient parallelisation.
\end{itemize}

\textbf{Future work.}  Planned extensions include: (i) MI-based feature
selection across a large number of variables using the pairwise matrix;
(ii) integration with standard neuroscience data formats (NWB, MNE-Python);
(iii) online estimation for streaming data; (iv) GPU-native multiprocessing.

%% ============================================================================
\section*{Acknowledgements}

\textit{I want to thank me for believing in me, I want to thank me for
doing all this hard work.  I wanna thank me for having no days off.
I wanna thank me for never quitting.  I wanna thank me for always being
a giver and trying to give more than I receive.  I wanna thank me for
trying to do more right than wrong.  I wanna thank me for being me at
all times.}
\hfill —Snoop Dogg

%% ---- Bibliography -----------------------------------------------------------
\bibliographystyle{plainnat}
\begin{thebibliography}{99}

\bibitem[Abdelaleem et al.(2025)]{abdelaleem2025accurate}
Abdelaleem, E., et al.\
``Accurate Estimation of Mutual Information in High Dimensional Data.''
\emph{arXiv:2506.00330}, 2025.

\bibitem[Abdelaleem et al.(2024)]{abdelaleem2024precision}
Abdelaleem, E., et al.\
``An information theoretic method to resolve millisecond-scale spike timing
precision in a comprehensive motor program.''
\emph{Proceedings of the National Academy of Sciences}, 2024.

\bibitem[Belghazi et al.(2018)]{belghazi2018mine}
Belghazi, M.\ I., et al.\
``Mutual Information Neural Estimation.''
\emph{Proceedings of the 35th ICML}, PMLR, 2018.

\bibitem[Borst and Theunissen(1999)]{borst1999information}
Borst, A.\ and Theunissen, F.\ E.\
``Information theory and neural coding.''
\emph{Nature Neuroscience}, 2(11):947--957, 1999.

\bibitem[Gulati et al.(2026)]{gulati2026mutual}
Gulati, T., Abdelaleem, E., et al.\
``Mutual Information Reveals the Dimensionality of Shared Information in
Neural Populations.''
\emph{arXiv:2602.08105}, 2026.

\bibitem[Kraskov et al.(2004)]{kraskov2004estimating}
Kraskov, A., St{\"o}gbauer, H., and Grassberger, P.\
``Estimating mutual information.''
\emph{Physical Review E}, 69(6):066138, 2004.

\bibitem[Oord et al.(2018)]{oord2018representation}
Oord, A.\ v.\ d., Li, Y., and Vinyals, O.\
``Representation Learning with Contrastive Predictive Coding.''
\emph{arXiv:1807.03748}, 2018.

\bibitem[Poole et al.(2019)]{poole2019variational}
Poole, B., Ozair, S., Oord, A.\ v.\ d., Alemi, A.\ A., and Tucker, G.\
``On Variational Bounds of Mutual Information.''
\emph{Proceedings of the 36th ICML}, PMLR, 2019.

\bibitem[Quian Quiroga and Panzeri(2009)]{quian2009measuring}
Quian Quiroga, R.\ and Panzeri, S.\
``Extracting information from neuronal populations: information theory and
decoding approaches.''
\emph{Nature Reviews Neuroscience}, 10(3):173--185, 2009.

\bibitem[Schreiber(2000)]{schreiber2000measuring}
Schreiber, T.\
``Measuring information transfer.''
\emph{Physical Review Letters}, 85(2):461--464, 2000.

\bibitem[Song and Ermon(2020)]{song2020understanding}
Song, J.\ and Ermon, S.\
``Understanding the Limitations of Variational Mutual Information Estimators.''
\emph{ICLR 2020}.

\bibitem[Timme and Lapish(2018)]{timme2018tutorial}
Timme, N.\ M.\ and Lapish, C.\
``A Tutorial for Information Theory in Neuroscience.''
\emph{eNeuro}, 5(3), 2018.

\end{thebibliography}

%% ============================================================================
\appendix

\section{Complete \texttt{nmi.run()} Parameter Reference}
\label{app:run_params}

\subsection{Top-level Parameters}

The following parameters are accepted by \api{nmi.run()} regardless of mode.

\begin{center}
\small
\begin{tabular}{lllp{7cm}}
\toprule
\textbf{Parameter} & \textbf{Type} & \textbf{Default} & \textbf{Description} \\
\midrule
\api{x\_data} & array-like & — & Input data for variable $X$. \\
\api{y\_data} & array-like & \api{None} & Input data for variable $Y$. \\
\api{mode} & str & \api{'estimate'} & Analysis mode (one of 9). \\
\api{base\_params} & dict & \api{\{\}} & Architecture + training params. \\
\api{processor\_type\_x} & str/None & \api{None} & \api{'continuous'}/\api{'spike'}/\api{'categorical'}/\api{None}. \\
\api{processor\_params\_x} & dict & \api{None} & Kwargs for the X processor. \\
\api{processor\_type\_y} & str/None & \api{None} & Same for $Y$. \\
\api{processor\_params\_y} & dict & \api{None} & Kwargs for the Y processor. \\
\api{output\_units} & str & \api{'nats'} & \api{'nats'} or \api{'bits'} ($\div\ln 2$). \\
\api{n\_workers} & int & 1 & Parallel processes. \\
\api{random\_seed} & int & \api{None} & NumPy/PyTorch seed. \\
\api{split\_mode} & str & \api{'blocked'} & \api{'blocked'} or \api{'random'}. \\
\api{train\_indices} & np.ndarray & \api{None} & Manual train indices. \\
\api{test\_indices} & np.ndarray & \api{None} & Manual test indices. \\
\api{permutation\_test} & bool & \api{False} & Run permutation null test. \\
\api{n\_permutations} & int & 100 & Permutations for null. \\
\api{bidirectional\_te} & bool & \api{False} & Compute $\mathrm{TE}(Y\to X)$ too. \\
\api{z\_data} & array-like & \api{None} & Conditioning variable for conditional MI. \\
\api{z\_time} & np.ndarray & \api{None} & Time vector for $Z$ dataset. \\
\api{z\_processor\_type} & str & \api{None} & Processor type for $Z$. \\
\api{z\_processor\_params} & dict & \api{None} & Processor params for $Z$. \\
\bottomrule
\end{tabular}
\end{center}

\subsection{\texttt{base\_params} Reference}

All model architecture and training parameters are passed via \api{base\_params}.

\begin{center}
\small
\begin{tabular}{lllp{6.5cm}}
\toprule
\textbf{Key} & \textbf{Type} & \textbf{Default} & \textbf{Description} \\
\midrule
\multicolumn{4}{l}{\textit{Architecture}} \\
\api{critic\_type} & str & \api{'separable'} & \api{'separable'}/\api{'hybrid'}/\api{'concat'}. \\
\api{embedding\_model} & str & \api{'mlp'} & \api{'mlp'}/\api{'cnn'}/\api{'gru'}/\api{'lstm'}/\api{'tcn'}/\api{'transformer'}. \\
\api{hidden\_dim} & int & 128 & Width of hidden layers. \\
\api{embedding\_dim} & int & 64 & Bottleneck dimension $d_z$. \\
\api{n\_layers} & int & 2 & Depth of the embedding network. \\
\api{use\_variational} & bool & \api{False} & Use variational embeddings. \\
\api{beta} & float & 1024.0 & Weight on MI in variational loss. \\
\api{shared\_encoder} & bool & \api{False} & Share embedding weights for $X$/$Y$. \\
\api{dropout} & float & 0.0 & Dropout probability. \\
\api{norm\_layer} & str & \api{None} & \api{'layer'}/\api{'batch'}/\api{None}. \\
\api{use\_spectral\_norm} & bool & \api{True} & Spectral normalisation on hidden layers. \\
\api{kernel\_size} & int & 7 & CNN/TCN kernel size (odd). \\
\api{nhead} & int & 4 & Transformer attention heads. \\
\api{custom\_critic} & nn.Module & \api{None} & Pre-built critic (bypasses build step). \\
\api{custom\_embedding\_cls} & class & \api{None} & Custom embedding class. \\
\midrule
\multicolumn{4}{l}{\textit{Training schedule}} \\
\api{estimator\_name} & str & \api{'infonce'} & \api{'infonce'}/\api{'smile'}/\api{'tuba'}/\api{'nwj'}/\api{'js\_fgan'}. \\
\api{estimator\_params} & dict & \api{None} & E.g.\ \api{\{'clip': 5.0\}} for SMILE. \\
\api{n\_epochs} & int & 50 & Maximum training epochs. \\
\api{batch\_size} & int & 256 & Samples per batch. \\
\api{learning\_rate} & float & 0.001 & Optimiser learning rate. \\
\api{optimizer} & str & \api{'adam'} & Optimiser name or class. \\
\api{optimizer\_params} & dict & \api{\{\}} & Extra kwargs for the optimiser. \\
\api{patience} & int & 10 & Early-stopping patience. \\
\api{smoothing\_sigma} & float & 1.0 & Gaussian smoothing $\sigma$. \\
\api{median\_window} & int & 5 & Median filter window. \\
\api{train\_fraction} & float & 0.9 & Train/test split fraction. \\
\api{n\_test\_blocks} & int & 5 & Blocks in blocked split. \\
\api{split\_gap\_fraction} & float & 0.5 & Gap fraction around test blocks. \\
\api{max\_eval\_samples} & int & 5000 & Max samples per MI evaluation. \\
\api{scheduler} & str & \api{None} & LR scheduler name or class. \\
\api{scheduler\_params} & dict & \api{\{\}} & Kwargs for the scheduler. \\
\api{gradient\_clip\_val} & float & \api{None} & Max gradient norm; \api{None} disables. \\
\api{spectral\_whitening} & str & \api{'std'} & \api{'std'}/\api{'zca'}/\api{None}. \\
\api{device} & str & \api{'auto'} & \api{'auto'}/\api{'cpu'}/\api{'cuda'}/\api{'mps'}. \\
\api{verbose} & bool & \api{False} & Print training details. \\
\api{show\_progress} & bool & \api{True} & Show tqdm progress bars. \\
\midrule
\multicolumn{4}{l}{\textit{Rigorous mode only}} \\
\api{delta\_threshold} & float & 0.10 & Quadratic curvature threshold. \\
\api{min\_gamma\_points} & int & 5 & Minimum linear-region points. \\
\api{confidence\_level} & float & 0.68 & CI level ($\approx 1\sigma$). \\
\api{min\_reliable\_samples} & int & 1000 & Minimum chunk size for warning. \\
\bottomrule
\end{tabular}
\end{center}

\subsection{Mode-Specific Keyword Arguments}

\begin{center}
\small
\begin{tabular}{lllp{7cm}}
\toprule
\textbf{Mode} & \textbf{Parameter} & \textbf{Default} & \textbf{Description} \\
\midrule
\api{sweep} & \api{sweep\_grid} & \api{\{\}} & Dict of param-name $\to$ list. \\
\api{lag} & \api{lag\_range} & — & Range of integer lags. \\
         & \api{equalize\_n} & \api{False} & Equalise sample count across lags. \\
\api{precision} & \api{tau\_grid} & — & List of corruption levels. \\
                & \api{corrupt\_target} & \api{'x'} & \api{'x'}/\api{'y'}/\api{'both'}. \\
                & \api{corruption\_method} & \api{'rounding'} & \api{'rounding'}/\api{'noise'}. \\
                & \api{n\_noise\_samples} & 50 & Noise realisations per $\tau$. \\
                & \api{threshold\_ratio} & 0.9 & Fraction of baseline MI. \\
\api{rigorous} & \api{gamma\_range} & \api{range(1,11)} & Subsampling gammas. \\
\api{dimensionality} & \api{split\_method} & \api{'random'} & \api{'random'}/\api{'spatial'}/\api{'temporal'}. \\
                     & \api{n\_splits} & 5 & Independent repetitions. \\
                     & \api{spectral\_mode} & \api{'summary'} & \api{'summary'}/\api{'full'}. \\
\api{transfer} & \api{history\_window} & — & Past window length. \\
               & \api{prediction\_horizon} & 1 & Future window length. \\
               & \api{bidirectional\_te} & \api{False} & Compute $\mathrm{TE}(Y\to X)$. \\
\api{pairwise} & \api{pairs} & \api{None} & Explicit list of $(i,j)$ pairs. \\
\bottomrule
\end{tabular}
\end{center}

\section{Minimal Worked Examples}
\label{app:examples}

\subsection{Bias-Corrected Estimation (Rigorous Mode)}
\begin{lstlisting}
import neural_mi as nmi

x, y = nmi.generators.generate_nonlinear_from_latent(
    n_samples=2500, latent_dim=10, observed_dim=50, mi=3.0
)

result = nmi.run(
    x_data=x, y_data=y,
    mode='rigorous',
    base_params={
        'n_epochs': 100,
        'batch_size': 256,
        'embedding_model': 'mlp',
        'embedding_dim': 64,
        'patience': 15,
    },
    split_mode='random',   # IID data
    n_workers=4,
    random_seed=42,
)
print(f"Corrected MI: {result.mi_estimate:.3f}"
      f" +/- {result.details['mi_error']:.3f} nats")
result.plot()  # bias-correction fit
\end{lstlisting}

\subsection{Transfer Entropy}
\begin{lstlisting}
import neural_mi as nmi
import numpy as np

# Simulate X driving Y with a 10-sample lag
T, C = 5000, 1
x = np.random.randn(T, C).astype(np.float32)
y = np.roll(x, 10, axis=0) + 0.3 * np.random.randn(T, C).astype(np.float32)

result = nmi.run(
    x_data=x, y_data=y,
    mode='transfer',
    base_params={'n_epochs': 80, 'embedding_model': 'gru', 'hidden_dim': 64},
    history_window=20,
    prediction_horizon=1,
    bidirectional_te=True,
)
print(f"TE(X->Y): {result.details['te_xy']:.3f} nats")
print(f"TE(Y->X): {result.details['te_yx']:.3f} nats")
print(f"DI: {result.details['directionality_index']:+.3f}")
\end{lstlisting}

\subsection{Intrinsic Dimensionality}
\begin{lstlisting}
import neural_mi as nmi

# 100 neurons, 2000 windows, 50 time-steps per window
import torch
x_pop = torch.randn(2000, 100, 50)

result = nmi.run(
    x_data=x_pop,
    mode='dimensionality',
    base_params={
        'n_epochs': 100,
        'embedding_model': 'mlp',
        'embedding_dim': 64,   # over-parameterised bottleneck
    },
    split_method='random',
    n_splits=5,
    n_workers=4,
)
print(f"Intrinsic dimensionality: "
      f"{result.details['participation_ratio']:.2f} "
      f"+/- {result.details['participation_ratio_std']:.2f}")
\end{lstlisting}

\subsection{Spike-Timing Precision}
\begin{lstlisting}
import neural_mi as nmi

# spike_times_x, spike_times_y: lists of spike time arrays (seconds)
result = nmi.run(
    x_data=spike_times_x,
    y_data=spike_times_y,
    mode='precision',
    processor_type_x='spike',
    processor_params_x={'window_size': 0.1, 'sample_rate': 1000},
    processor_type_y='spike',
    processor_params_y={'window_size': 0.1, 'sample_rate': 1000},
    base_params={'n_epochs': 100, 'patience': 1000, 'batch_size': 512},
    tau_grid=[0.001, 0.002, 0.005, 0.010, 0.020, 0.050],
    threshold_ratio=[0.9, 0.75, 0.5],   # three thresholds at once
)
result.summary()
\end{lstlisting}

\subsection{All-to-All Connectivity Matrix}
\begin{lstlisting}
import neural_mi as nmi

# x_eeg: (n_windows, n_channels, window_size)
result = nmi.run(
    x_data=x_eeg,
    mode='pairwise',
    base_params={'n_epochs': 60, 'embedding_dim': 32},
    n_workers=8,
)
df = result.details['dataframe']
mi_matrix = result.details['mi_matrix']
print(df.sort_values('mi_mean', ascending=False).head(10))
\end{lstlisting}

\end{document}
